\section{Data and experimental setup}\label{sec:data}

\subsection{Controlled azimuth campaign (2026)}\label{subsec:campaign}
The primary validation dataset consists of measurements of GTI acquired at multiple documented orientations of a fixed tilted surface. In order to achieve this, an inclined Kipp and Zonen CMP6 pyranometer, with a tilt of 35 degrees, was mounted on a rotating graduated plate located at the rooftop of the Faculty of Engineering in Montevideo, Uruguay (Latitude: $ \SI{-34.92}{\degree}$ , Longitude: $ \SI{-56.17}{\degree}$, Elevation: ???).  Special care was taken when leveling the table where the plate is mounted as well as the inclined plane where the pyranometer is placed.

GHI simultaneous data is acquired with a horizontal Kipp and Zonen (MODEL?) pyranometer, both calibrated and with traceability to the WMO (DETALLES?) and regularly maintained. The data is collected at 1-min intervals.

True north ($\gamma = \SI{0}{\degree}$) was obtained with the help of the shadow of a vertical rod under clear sky conditions at the solar noon, as provided by the NASA solar calculator (WEBSITE?). 
By comparing the north oriented GTI data under clear sky conditions with simultaneous GHI data, it was established that the misalignment of the tilted sensor was minimal  \draftnote{Assess uncertainty from the shadow thickness in cm, and discuss the gamma =0 correct orientation}


\begin{table}[t]
\centering
\caption{Provisional inventory of the 2026 controlled experiment. Fill in from
measurement metadata; do not infer reference precision from the step size.}
\label{tab:campaign}
\begin{tabular}{llll}
\toprule
Period / date & $\beta_{\mathrm{ref}}$ ($^\circ$) & $\gamma_{\mathrm{ref}}$ ($^\circ$) & Valid days \\
\midrule
To be completed & -- & -- & -- \\
\bottomrule
\end{tabular}
\end{table}

\subsection{Reference orientation and its uncertainty}
\label{subsec:reference}

The geographical North–South reference was established independently of the irradiance measurements using the shadow of a vertical rod at the predicted solar noon. The procedure was performed in late May using a cylindrical rod approximately 1 m high, mounted on a leveled horizontal surface. Rod verticality was checked against its supporting base using a right-angle reference.
The resulting meridian line was marked directly on the supporting table and used to establish the zero-azimuth reference of the rotating platform. The platform featured angular graduations at 1° intervals, a bubble level, and mechanical fixation screws.
The uncertainty associated with this reference procedure was explicitly considered when assessing absolute azimuth errors. Additionally, relative angular displacements were evaluated independently of the common initial alignment error.


We distinguish between nominal azimuth ($\gamma_{\mathrm{nom}}$), corresponding to the reported or intended surface orientation, and reference azimuth ($\gamma_{\mathrm{ref}}$), established through an independent geometric procedure. In the controlled experimental campaign, reference orientations were determined using an astronomically established North–South direction and a graduated rotating platform. For the remote stations, only nominal orientations were generally available, and the estimated azimuths were therefore interpreted as potential orientation corrections rather than independently validated measurements.
\draftnote{Document the reference procedure (survey, calibrated instrument, alignment, fixture), repeatability and estimated uncertainty for tilt and azimuth. Avoid claiming absolute angular accuracy better than the reference allows.}

Esto es un tema a discutir - Agus?


\subsection{Quality control and dataset construction}
\label{subsec:datasetqc}
All quality-control procedures should be specified before evaluating their
impact on azimuth errors. In addition to physical plausibility screening, the
availability of suitable clear-sky intervals around the GTI maximum must be
reported together with the fraction of excluded samples.
\draftnote{Create an auditable QC log: raw count -> daytime -> clear sky -> fit window -> outlier rejection -> valid estimates, with per-day reasons. State whether masks use GTI alone or additional channels.}


\subsection{Independent field applications}\label{subsec:fielddata}
To assess operational usefulness, the method will also be applied to
historical GTI measurements in which the reported orientation may be
uncertain or time-dependent \citep{piccioli2022thesis}.

\begin{itemize}
	\item Estación TA, describir equipos, timespan, medidas disponibles, instrumentación y describir evidencia de que la misma se encuentra desviada un ángulo indeterminado
	\item Esto resulta en modelos de transposición cuya validación es errónea, con métricas que no coinciden con lo esperado según la literatura
	\item GTI en cielo claro muestra un máximo en $\omega^{\ast} \neq 0$
	\item Reportar cantidad de días claros evaluados, valores obtenidos, promedio de azimuts
	\item Mostrar métricas con el valor de azimut calculado y ver cómo los porcentajes se acercan a lo esperado y conocido en la literatura.
\end{itemize}

\draftnote{Select and describe stations from the thesis; add site identifiers, available periods, instrumentation, and any independent evidence of orientation changes. State clearly which cases are diagnoses rather than externally validated ground truth.}
